\section{The Pristine Theology}

I had a dream two nights ago. A voice was giving me great wisdom. It said it was a blessing to have access to such wisdom while I was sleeping and that I should share it. I could barely wait to wake up so I could write it all down. When I awoke, I grabbed a pen, but there was nothing there to write about. So I consulted the Akashic record, but it was all pixelated, and this is all I could make out.

Take it as a warning to remain awake when receiving revelations from the Divine Mind.

\paragraph{Synarchy}
There are a great many theories about the proper socio-political organization of societies. However, only those based in reality, viz., spiritual reality, are worth pursuing, and not the various factitious or artificial constructions. Specifically, there is a spiritual authority, a political power, and economic activity, hierarchically arranged. Within this arrangement, there are still very many possible variations. The word for this arrangement is Synarchy. Of course, authority must be intelligent and power just. There is no escape from this system. The alternate is an ersatz version of each level.

Some quick examples. In the Middle Ages, there was the Peace and Truce of God. This limited the political powers from tramping on the common people, otherwise they accept no limits. In Europe, there was the ideal of the gentleman. To the extent they themselves aspired to that ideal, the money powers limited themselves to fair play. Without it, they admit no limits to the pursuit of wealth.

\paragraph{The Perennial Philosophy}
It is false that the West has no authentic tradition, and \textbf{Rene Guenon} admits it in his study of Hindu Doctrine. This has gone under the name of \emph{prisca theologia} (pristine theology) or perennial philosophy. Following Ficino, \textbf{Valentin Tomberg} adheres to the same tradition as we indicated three years ago in \textit{Heart of the West}\footnote{\url{http://www.gornahoor.net/?p=1256}}. What Tomberg meant is that he is following the Hermetic tradition in the West.

Nevertheless, there are still those who believe that the most important question is whether or not Tomberg remained an Anthroposophist. Hence, they call Tomberg a ``Platonist", even though Plato was but one link in that chain, since that is a fundamental anthroposophic category.

\paragraph{Continuity of Tradition}
Guenon was not the first to deny the continuity of the Western tradition, since Ficino's timeline was debunked a few centuries ago through textual criticism. However, a Tradition is not simply a continuity of texts. In Ficino's view, the Pristine Theology derived from the Divine Mind, not what someone wrote previously.

Fortunately, Tomberg has given us a complete description of how a revelation becomes gnosis and finally a book. Textual criticism is irrelevant in this case; all that matters is the revelation and how well it is expressed. Of course, you would have to learn to think with your heart rather than with your head.

Now a theologian as sophisticated as \textbf{Vladimir Solovyov} claims that Christian dogma itself can be traced back to Egyptian Hermetism. Since he never wrote his promised defense of that idea, probably because it was so obvious to him, an intelligent man today may want to take on that task now.

\paragraph{Arcana and Symbols}
Although Guenon is of utmost importance for demonstrating the universal nature of symbols, Tomberg has given us an even deeper understanding of symbolism. Saying that a symbol ``is" really ``that", is a start, but it is quite insufficient. Otherwise, we would be in the position of the person who gave me, in a seminar on Meditations on the Tarot, the signs in an airport as examples of symbols. They are univocal, but there is no value to meditating on a no smoking sign or the sign in front of the men's restroom.

Guenon often mentioned the symbol of the two birds in the fig tree, one silent, the other active. One could easily relate that to the transcendent self, observing the activities of the ego. But that gets you nowhere. Instead, you should meditate carefully on the entire symbol, drawing out its full import. That makes it an Arcanum, not just a symbol.

\paragraph{The Center is Everywhere}
\textbf{Miguel Serrano} mentions this quote from \textbf{Carl Jung} from one of their meetings (although Jung spoke it in Latin):

\begin{quotex}
The Self is a circle whose center is everywhere and whose circumference is nowhere. 

\end{quotex}
Surprisingly, Serrano did not catch the allusion. This description was applied to God by \textbf{Alain de Lille} and, I believe, \textbf{Nicholas Cusanus}. So Jung is implying that the Self is God. Of course, this is no surprise, since Serrano then quotes Jung as saying that this Self is Christ, at least for Western man, since

\begin{quotex}
[Christ is] the archetype of the hero, representing man's highest aspiration. All this is very mysterious and at times frightening. 

\end{quotex}
Nevertheless, that is the path to the Self and Jesus ``meek and mild" will not get you there.

\paragraph{Theology of the Body}
I have a 700+ page book on the theology of the body by Pope, now Saint, \textbf{John Paul II}. I believe that a 700 page book on the theology of the soul would have been more appropriate, since there is very little knowledge of the soul today. Nevertheless, I gave the book a try. Wishing to read instead the executive summary, I discovered that there is a cottage industry offering classes on that book, for a price of course.

One of first things I found out is that it is not sinful for a man to have anal intercourse with his wife as long as he then ejaculates in her vagina. My first reaction was, why would anyone want to do that, but I realized people were paying money to learn such things. So I put the book down.

\paragraph{Buddhist Sex}
Many years ago, I was still under Guenon's influence on this issue and decided it was necessary to be initiated into a Tradition. I was then initiated into Tibetan Buddhism. So for many years I practiced Tibetan and Zen meditation, chanted the \emph{Heart Sutra}, stared at Thangkas, and so on. I am not recommending it, because the caliber is very low, even though I heard countless times how the aspiring Buddhists had passed beyond the religion of their ancestors that they had learned in childhood. Nevertheless, the parallels are there, as might be expected if the various traditions do indeed represent some common Tradition.

First of all, there are ``venial" and ``mortal" sins. The former can be fixed by an act of purification, but the latter requires a retaking of one's vows publicly. As an example of a mortal sin (one of just four), there is this:

\begin{quotex}
engaging in sexual misconduct, that is, engaging in sexual activities with another's husband or wife, or using the mouth or anus of one's own husband or wife. 

\end{quotex}
If you read quaint old books, you will recognize these activities as adultery and sodomy. So this is a stricter standard than even JP II and may explain my prudishness. This is the one sentence version of the theology of the body, although if you look, and not even very hard, you will find a lot of words to mitigate the force of that transgression.

For the record, the other transgressions are murder, stealing something of value, and claiming to be enlightened when you are not. No commentary is necessary.



\flrightit{Posted on 2014-06-19 by Cologero }

\begin{center}* * *\end{center}

\begin{footnotesize}\begin{sffamily}



\texttt{Kimiwoux on 2014-06-19 at 13:16 said: }

Are there examples of Synarchy in older societies? I don't believe the US is a Synarchy. I guess that's an understatement. Maybe the period of the Knights Templar?

I'm starting to learn through a gnostic tradition that these sexual codes were important. It's changed the way I view sex and life. This was a really good post. Glad I squadalahed over here.


\hfill

\texttt{Tomas Diaz on 2014-06-19 at 16:34 said: }

I would hope that one not understand the book from the cottage industry. Men like Christopher West are, more slowly than some would like, being discredited as an interpreter of the Pope's work. John Paul's own title for the catechesis being ``Man and Woman He Created Them" – it is a theology of sexuality, understood in the most spiritual sense, and the ethics that West and others seek to promote out of it are contrary to the Pope's intention.

John Paul does not discuss the merits or lack thereof of sodomy. He assumed it be obvious sin.


\hfill

\texttt{Cologero on 2014-06-19 at 17:29 said: }

Thanks for pointing out, Mr. Diaz, that the original work cannot be judged by the commentaries. I was being a little satirical and meant no disrepect to Pope John Paul II.


\hfill

\texttt{Synodius on 2014-06-19 at 20:17 said: }

Some historical background on JPII's ideas in his TOB

\url{http://www.traditioninaction.org/HotTopics/f024ht\_Hildebrand\_Engel.htm}

Primary and secondary ends of marriage and hierarchy between these is an interesting topic and I feel like the traditional focus on procreation needed to be balanced, it is wierd to give primacy to lower, physical reasons.

Dalai Lama seems to go with the flow these days…

\url{http://abcnews.go.com/Politics/dalai-lama-says-he-supports-gay-marriage/blogEntry?id=22834522}


\hfill

\texttt{Cologero on 2014-06-20 at 00:23 said: }

I agree, Synodius, that some balance is necessary. Marriage itself is more than procreation. Fr. Ripperger's sermons on Marriage\footnote{\url{http://www.sensustraditionis.org/webaudio/Marriage/Marriage.mp3}} and Women and the Natural Order\footnote{\url{http://www.sensustraditionis.org/webaudio/Sermons/KC/Disk\%201/Women.mp3}} may be of interest. There, he is claiming that the primary end of marriage is the salvation of each others' souls. That is interesting from the exoteric point of view.

The esoteric view builds on that notion through the ideal of courtly love or Mouravieff's polar beings. However, those relations do not necessarily involve physical sexual relations. Miguel Serrano's Test\footnote{\url{http://www.gornahoor.net/?p=70}} is instructive in that regard.


\hfill

\texttt{CaseyAnn on 2014-06-20 at 10:50 said: }

I'm glad to be informed of this. And thank you, Synodius, for linking to Tradition in Action as a counterpoint to the research I did after reading this last night. Alice von Hildebrand critiques West's work very seriously\footnote{\url{http://www.catholicnewsagency.com/document.php?n=999}}. But (at least I think) one should always check TIA. 

``is Christ, at least for Western man" …

Cologero, I'm sorry to say that I will be taking Gornahoor off my blogroll. At this point I'm now in an area where I can go to the Tridentine Mass and am much more conscious of the Tradition of the Church, servitude to Our Lady and so forth. While I acknowledge the noble intentions here (and have personally defended them against fellow traditionalist Catholics) and the efficacy of the pre-Christian religions, I believe one is bound to Christ explicitly when one hears the true religion proclaimed. 

Que Dieu vous bénisse.


\hfill

\texttt{Cologero on 2014-06-23 at 10:57 said: }

Thanks, Casey Ann, for your support and defense. You are free to make any choices you prefer. I have mentioned on more than one occasion that Gornahoor is optional and is not necessary. Nevertheless, some people respond better to it than more rigid exoteric approaches; it is based on valid traditions and the sources are documented.

I wish to remind careful readers that quoting an author is not tantamount to endorsing all their opinions. Certain messages are intended for some readers rather than others; everything must be read as a whole. For example, the Jung quote was intended to evoke ``putting on the mind of Christ". That is the true religion.

I should also remind readers not to confuse specific devotions with the whole of Tradition. For example, ``servitude to Our Lady" is neither a dogma nor a Tradition (in our sense), but rather a devotion used by certain orders. It is not binding on all.


\hfill

\texttt{CaseyAnn on 2014-06-23 at 12:57 said: }

Thank you, Cologero. I should have been more careful in using ``Tradition" here; I actually meant that specifically in reference to the Church, which would include devotion to Our Lady. I do think this St. Louis de Montfort's take on this is very important, however, so I'll leave you all with it: ``He who has not Mary for his Mother has not God for his Father."


\end{sffamily}\end{footnotesize}
